\begin{theorem}[\texorpdfstring{$\mathbb Z_2^2$}{Z2^2} 指纹的中心幂等元实现与扇区理想的显式投影]\label{thm:fold-groupoid-z2x2-central-idempotents}
令 \(\mathcal{R}_m=\Omega_m\times_{X_m}\Omega_m\) 为定义 \ref{def:fold-equivalence-groupoid-algebra} 的折叠等价关系群胚，并把其群胚代数记为
\[
A_{m,1}:=\CC[\mathcal{R}_m]\ \cong\ \bigoplus_{x\in X_m}M_{d_m(x)}(\CC)
\]
（命题 \ref{prop:fold-groupoid-wedderburn}）。
对每个 \(x\in X_m\)，令 \(\widetilde P_x\in Z(A_{m,1})\) 为对应简单块的极小中心幂等元，使得
\(
\widetilde P_x A_{m,1}\cong M_{d_m(x)}(\CC)
\)。
并令扫描—反演取向指纹
\(
\chi(x)=(\operatorname{sgn}(c_x),\operatorname{sgn}(\iota_x))\in\{\pm1\}^2
\)
如命题 \ref{prop:fold-groupoid-z2x2-sector-by-orientation}。
\par
定义两个中心自伴酉元（中心对合）
\[
\varepsilon_{\mathrm{scan}}:=\sum_{x\in X_m}\operatorname{sgn}(c_x)\,\widetilde P_x,
\qquad
\varepsilon_{\mathrm{rev}}:=\sum_{x\in X_m}\operatorname{sgn}(\iota_x)\,\widetilde P_x.
\]
则 \(\varepsilon_{\mathrm{scan}}^2=\varepsilon_{\mathrm{rev}}^2=1\)，且 \(\varepsilon_{\mathrm{scan}}\varepsilon_{\mathrm{rev}}=\varepsilon_{\mathrm{rev}}\varepsilon_{\mathrm{scan}}\)。
对任意 \((\alpha,\beta)\in\{\pm1\}^2\)，定义
\[
e_{\alpha,\beta}
:=
\frac14\bigl(1+\alpha\,\varepsilon_{\mathrm{scan}}\bigr)\bigl(1+\beta\,\varepsilon_{\mathrm{rev}}\bigr)
\ \in\ Z(A_{m,1}).
\]
则 \(\{e_{\alpha,\beta}\}_{(\alpha,\beta)\in\{\pm1\}^2}\) 构成一组两两正交、和为 \(1\) 的中心幂等元，并且对应的扇区理想满足
\[
e_{\alpha,\beta}A_{m,1}
\ \cong\
\bigoplus_{\chi(x)=(\alpha,\beta)}M_{d_m(x)}(\CC).
\]
从而 \(\mathbb Z_2^2\) 扇区分解不仅是索引集合的分解，而且可由 \((\varepsilon_{\mathrm{scan}},\varepsilon_{\mathrm{rev}})\) 的联合谱分解在代数内部显式实现。
此外，由 \(\varepsilon_{\mathrm{scan}},\varepsilon_{\mathrm{rev}}\) 生成的读出子代数
\[
\mathcal A_\chi:=C^\ast(\varepsilon_{\mathrm{scan}},\varepsilon_{\mathrm{rev}})
\]
为四维交换半单 \(C^\ast\)-代数，且典范同构于 \(\CC[\mathbb Z_2^2]\cong \CC^4\)。
\leanverified{wedderburn\_dim\_m6}
\leanverified{fiber\_histogram\_m6}
\leanverified{projectorVal}
\leanverified{projectorVal\_partition\_of\_signs}
\leanverified{projectorVal\_case\_split}
\leanverified{epn\_idempotent}
\leanverified{enp\_idempotent}
\leanverified{enn\_idempotent}
\leanverified{projectorVal\_eq\_one\_iff}
\leanverified{projectorVal\_zero\_iff\_ne}
\leanverified{projectorVal\_sum\_eq\_one\_of\_cases}
\leanverified{projectorVal\_eq\_one\_iff\_other\_three\_zero}
\leanverified{paper\_projectorVal\_eq\_one\_iff\_other\_three\_zero}
\leanverified{projectorVal\_ne\_neg\_one}
\leanverified{projectorVal\_zero\_iff\_other\_exists\_one}
\leanverified{projectorVal\_hadamard\_sector\_powersum}
\leanverified{central\_idempotents\_pairwise\_orthogonal}
\leanverified{central\_idempotents\_all\_idempotent}
\leanverified{paper\_wedderburn\_central\_idempotent\_package}
\leanverified{epp\_mul\_epn}
\leanverified{epp\_mul\_enp}
\leanverified{epp\_mul\_enn}
\leanverified{epn\_mul\_enp}
\leanverified{epn\_mul\_enn}
\leanverified{enp\_mul\_enn}
\leanverified{paper\_central\_idempotents\_complete\_system}
\leanverified{paper\_fold\_groupoid\_z2x2\_central\_idempotents}
\leanverified{paper\_sector\_dimension\_sum\_m6}
\leanverified{wedderburn\_dim\_ratio\_m7\_exact}
\leanverified{wedderburn\_dim\_ratio\_m6\_not\_exact}
\leanverified{projectorVal\_at\_same\_signs}
\leanverified{paper\_projectorVal\_sign\_match\_complete}
\leanverified{central\_idempotents\_at\_one\_one}
\leanverified{central\_idempotents\_at\_neg\_one\_neg\_one}
\leanverified{central\_idempotents\_at\_one\_neg\_one}
\leanverified{central\_idempotents\_at\_neg\_one\_one}
\leanverified{paper\_central\_idempotents\_four\_corner\_witness}
\leanverified{projectorVal\_at\_one\_one\_one\_one}
\leanverified{projectorVal\_at\_mismatch\_zero}
\leanverified{projectorVal\_at\_neg\_one\_one\_neg\_one\_one}
\leanverified{projectorVal\_at\_all\_neg\_one}
\leanverified{paper\_projectorVal\_corner\_table}
\end{theorem}

\begin{proof}
由 Wedderburn 分解，中心 \(Z(A_{m,1})\) 由 \(\{\widetilde P_x\}\) 线性张成，且
\(\widetilde P_x\widetilde P_y=\delta_{xy}\widetilde P_x\)、\(\sum_x\widetilde P_x=1\)。
因此 \(\varepsilon_{\mathrm{scan}},\varepsilon_{\mathrm{rev}}\) 均为中心元，并在每个简单块 \(M_{d_m(x)}(\CC)\) 上取值为标量 \(\operatorname{sgn}(c_x)\) 与 \(\operatorname{sgn}(\iota_x)\)，从而自伴且平方为 \(1\)，并且彼此可交换。
\par
由可交换性直接计算得 \(e_{\alpha,\beta}^2=e_{\alpha,\beta}\)，且当 \((\alpha,\beta)\neq(\alpha',\beta')\) 时有 \(e_{\alpha,\beta}e_{\alpha',\beta'}=0\)，并满足 \(\sum_{\alpha,\beta}e_{\alpha,\beta}=1\)。
在块 \(x\) 上，
\[
e_{\alpha,\beta}\Big|_x
=
\frac14\bigl(1+\alpha\,\operatorname{sgn}(c_x)\bigr)\bigl(1+\beta\,\operatorname{sgn}(\iota_x)\bigr),
\]
该标量当且仅当 \(\chi(x)=(\alpha,\beta)\) 时为 \(1\)，否则为 \(0\)。
因此 \(e_{\alpha,\beta}\) 恰好在扇区 \(\chi(x)=(\alpha,\beta)\) 的块上取恒等，而在其余块上为零，从而
\(
e_{\alpha,\beta}A_{m,1}\cong \bigoplus_{\chi(x)=(\alpha,\beta)}M_{d_m(x)}(\CC)
\)。
\par
最后，\(\mathcal A_\chi\) 由两个可交换对合生成，因而与 \(\CC\)-代数 \(\CC[\mathbb Z_2^2]\) 的生成关系一致；其极小中心幂等元正是 \(\{e_{\alpha,\beta}\}\)，故 \(\mathcal A_\chi\cong \CC^4\)。证毕。
\end{proof}

\begin{theorem}[\(\chi\) 的联合谱测度与推前分布一致：带符号碰撞矩的迹公式]\label{thm:fold-groupoid-z2x2-joint-spectral-measure}
在 \(A_{m,1}\cong\bigoplus_{x\in X_m}M_{d_m(x)}(\CC)\) 上取归一化迹
\[
\tau_{m,1}(A):=2^{-m}\sum_{x\in X_m}\Tr(A_x),
\]
其中 \(A_x\in M_{d_m(x)}(\CC)\) 为第 \(x\) 块分量。
则对每个极小中心投影 \(\widetilde P_x\) 有
\[
\tau_{m,1}(\widetilde P_x)=\frac{d_m(x)}{2^m}=: \pi_m(x),
\]
即 \(\pi_m\) 为均匀微态基线在折叠 \(\Fold_m\) 下的推前分布。
\par
令 \((\varepsilon_{\mathrm{scan}},\varepsilon_{\mathrm{rev}})\) 如定理 \ref{thm:fold-groupoid-z2x2-central-idempotents}。
则对任意函数 \(f:\{\pm1\}^2\to\CC\)，有恒等式
\[
\tau_{m,1}\!\bigl(f(\varepsilon_{\mathrm{scan}},\varepsilon_{\mathrm{rev}})\bigr)
=\sum_{x\in X_m}\pi_m(x)\,f\!\bigl(\chi(x)\bigr).
\]
从而在 \(\tau_{m,1}\) 下，\((\varepsilon_{\mathrm{scan}},\varepsilon_{\mathrm{rev}})\) 的联合谱测度与 \(\chi\) 在 \(\pi_m\) 下的推前测度 \(\chi_\ast\pi_m\) 严格一致。
\par
进一步，令
\[
\mathrm{IndW}(E_{m,1})=\sum_{x\in X_m} d_m(x)\,\widetilde P_x\ \in\ Z(A_{m,1})
\]
为分块退极化条件期望 \(E_{m,1}\) 的 Watatani 指标元。
则对任意整数 \(q\ge 0\) 与 \(s,t\in\{0,1\}\) 有闭式
\[
\tau_{m,1}\!\bigl(\varepsilon_{\mathrm{scan}}^{\,s}\varepsilon_{\mathrm{rev}}^{\,t}\,\mathrm{IndW}(E_{m,1})^{\,q}\bigr)
=
\frac{1}{2^m}\sum_{x\in X_m}\operatorname{sgn}(c_x)^{s}\operatorname{sgn}(\iota_x)^{t}\,d_m(x)^{q+1}.
\]
\end{theorem}
\leanverified{paper\_fold\_groupoid\_z2x2\_joint\_spectral\_measure}

\begin{proof}
由中心分解与 \(\varepsilon_{\mathrm{scan}},\varepsilon_{\mathrm{rev}}\) 在块 \(x\) 上取标量 \(\chi(x)\)，可写
\[
f(\varepsilon_{\mathrm{scan}},\varepsilon_{\mathrm{rev}})
=\sum_{x\in X_m} f(\chi(x))\,\widetilde P_x.
\]
取迹并用 \(\tau_{m,1}(\widetilde P_x)=2^{-m}\Tr(I_{d_m(x)})=d_m(x)/2^m\) 即得第一式。
\par
第二式同理：\(\mathrm{IndW}(E_{m,1})\) 在块 \(x\) 上取标量 \(d_m(x)\)，故
\[
\varepsilon_{\mathrm{scan}}^{\,s}\varepsilon_{\mathrm{rev}}^{\,t}\,\mathrm{IndW}(E_{m,1})^{\,q}
=\sum_{x\in X_m}\operatorname{sgn}(c_x)^s\operatorname{sgn}(\iota_x)^t\,d_m(x)^q\,\widetilde P_x,
\]
再取迹得到
\[
2^{-m}\sum_{x\in X_m}\operatorname{sgn}(c_x)^s\operatorname{sgn}(\iota_x)^t\,d_m(x)^q\,\Tr(I_{d_m(x)})
=\frac{1}{2^m}\sum_{x\in X_m}\operatorname{sgn}(c_x)^s\operatorname{sgn}(\iota_x)^t\,d_m(x)^{q+1}.
\]
证毕。
\end{proof}

\begin{proposition}[极大矩阵理想的 \texorpdfstring{$\chi$}{chi}-同质性：最大块仅落在单一扇区]\label{prop:fold-groupoid-maxblock-chi-homogeneity}
令 \(\mathcal{M}_m:=\{x\in X_m:\ d_m(x)=D_m\}\)，并令极大块理想
\[
I_m:=\bigoplus_{x\in\mathcal{M}_m}M_{D_m}(\CC)\ \subset\ A_{m,1}\cong\bigoplus_{x\in X_m}M_{d_m(x)}(\CC).
\]
设 \(\chi_m\in\{\pm1\}^2\) 为定理 \ref{thm:pom-maxfiber-orientation-charge-periodic} 给出的最大纤维族取向指纹常值，并记 \(e_{\chi_m}:=e_{\alpha,\beta}\)（其中 \(\chi_m=(\alpha,\beta)\)）为定理 \ref{thm:fold-groupoid-z2x2-central-idempotents} 的中心幂等元。
则有包含关系
\[
\boxed{\;I_m\subseteq e_{\chi_m}A_{m,1}.\;}
\]
特别地，商代数 \(A_{m,1}/e_{\chi_m}A_{m,1}\) 的 Wedderburn 直和中不含任何维数为 \(D_m\) 的矩阵块。
\leanverified{paper\_fold\_groupoid\_maxblock\_chi\_homogeneity}
\end{proposition}
\begin{proof}
由定理 \ref{thm:pom-maxfiber-orientation-charge-periodic}，对任意 \(x\in\mathcal{M}_m\) 有 \(\chi(x)=\chi_m\)。
而由定理 \ref{thm:fold-groupoid-z2x2-central-idempotents}，中心投影 \(e_{\chi_m}\) 在块 \(x\) 上取值为 \(1\) 当且仅当 \(\chi(x)=\chi_m\)，否则为 \(0\)。
因此对每个 \(x\in\mathcal{M}_m\) 有 \(\widetilde P_x\le e_{\chi_m}\)，从而
\(
M_{D_m}(\CC)\cong \widetilde P_xA_{m,1}\subseteq e_{\chi_m}A_{m,1}
\)；
对 \(x\in\mathcal{M}_m\) 求直和即得 \(I_m\subseteq e_{\chi_m}A_{m,1}\)。
最后一句为该包含关系在 Wedderburn 直和中的直接释义。证毕。
\end{proof}

\begin{theorem}[最大数据处理损失的支撑刚性与 \texorpdfstring{$\chi$}{chi}-异常不变量]\label{thm:fold-max-data-processing-saturation-chi}
令 \(E_m\) 为折叠 \(\Fold_m:\Omega_m\to X_m\) 诱导的条件期望/粗粒化通道（沿纤维均匀平均），并定义最大相对熵落差
\[
\Delta_m:=\sup_{\rho,\sigma}\Bigl(D(\rho\Vert\sigma)-D(\rho\circ E_m\Vert\sigma\circ E_m)\Bigr),
\]
其中 \(\rho,\sigma\) 取遍 \(\Omega_m\) 上的概率分布（等价地取遍交换代数 \(A_m=\CC^{\Omega_m}\) 上的态）。
则：
\begin{enumerate}
\normalfont
  \item \(\Delta_m=\log D_m\)（命题 \ref{prop:fold-condexp-index-maxfiber}）。
  \item 若 \((\rho,\sigma)\) 达到上确界，则 \((\Fold_m)_\ast\rho\) 的支撑包含于 \(\mathcal{M}_m:=\{x\in X_m:\ d_m(x)=D_m\}\)；等价地，
  \[
  (\Fold_m)_\ast\rho(x)>0\ \Longrightarrow\ d_m(x)=D_m.
  \]
  \item 因而任取达到上确界的极值对 \((\rho,\sigma)\)，其宏观支撑落点 \(x^\star\in \mathrm{supp}((\Fold_m)_\ast\rho)\) 满足
  \[
  \chi(x^\star)=\chi_m,
  \]
  其中 \(\chi_m\) 为定理 \ref{thm:pom-maxfiber-orientation-charge-periodic} 的最大纤维族取向指纹常值。特别地，当 \(m\) 为奇数且处于稳定阈值区间时，有 \(\operatorname{sgn}(\iota_{x^\star})=+1\)。
\end{enumerate}
\end{theorem}
\leanverified{paper\_fold\_max\_data\_processing\_saturation\_chi}
\begin{proof}
第一条由命题 \ref{prop:fold-condexp-index-maxfiber} 给出。
\par
对第二条，令 \(X:=\Fold_m(\omega)\) 为宏观变量，并记 \(\rho_X:=(\Fold_m)_\ast\rho\)、\(\sigma_X:=(\Fold_m)_\ast\sigma\)。
对每个 \(x\in X_m\) 令 \(\rho(\cdot\mid x)\)、\(\sigma(\cdot\mid x)\) 为在纤维 \(\Fold_m^{-1}(x)\) 上的条件分布（在 \(\rho_X(x),\sigma_X(x)>0\) 处定义）。
则相对熵在粗粒化下满足分解恒等式
\[
D(\rho\Vert\sigma)-D(\rho_X\Vert\sigma_X)
=\sum_{x\in X_m}\rho_X(x)\,D\!\bigl(\rho(\cdot\mid x)\,\Vert\,\sigma(\cdot\mid x)\bigr).
\]
并且对每个 \(x\) 有上界
\(
D(\rho(\cdot\mid x)\Vert\sigma(\cdot\mid x))\le \log d_m(x)
\)
（因两分布皆支撑于大小为 \(d_m(x)\) 的有限集合）。
于是得到
\[
\Delta_m\le \sup_{\rho}\ \sum_{x}\rho_X(x)\,\log d_m(x)\ \le\ \log D_m.
\]
若 \((\rho,\sigma)\) 达到 \(\Delta_m=\log D_m\)，则上式中两处不等式必须同时取等。
特别地，第二个不等式的等号强制 \(\rho_X\) 的支撑仅落在满足 \(\log d_m(x)=\log D_m\) 的点上，即 \(d_m(x)=D_m\)，得到第二条。
\par
第三条由第二条与定理 \ref{thm:pom-maxfiber-orientation-charge-periodic} 的常值性推出：对任意 \(x^\star\in\mathrm{supp}(\rho_X)\subseteq\mathcal{M}_m\) 有 \(\chi(x^\star)=\chi_m\)。
奇数情形的第二分量恒为 \(+1\) 亦由该定理给出。证毕。
\end{proof}

\leanverified{paper\_fold\_indw\_high\_q\_extract\_maxfiber\_and\_chi}
\begin{corollary}[Watatani 指标元高阶矩的 \texorpdfstring{$q\to\infty$}{q->infty} 极限：最大纤维规模与 \texorpdfstring{$\chi_m$}{chi\_m} 指纹读出]\label{cor:fold-indw-high-q-extract-maxfiber-and-chi}
固定分辨率 \(m\)，令
\[
\mathrm{IndW}(E_{m,1})=\sum_{x\in X_m}d_m(x)\,\widetilde P_x\ \in\ Z(A_{m,1})
\]
如定理 \ref{thm:fold-groupoid-z2x2-joint-spectral-measure}。
记
\[
D_m:=\max_{x\in X_m}d_m(x),\qquad
\mathcal M_m:=\{x\in X_m:\ d_m(x)=D_m\},
\qquad
\chi_m=(\alpha_m,\beta_m)\in\{\pm1\}^2
\]
其中 \(\chi_m\) 由定理 \ref{thm:pom-maxfiber-orientation-charge-periodic} 给出，且由命题 \ref{prop:fold-groupoid-maxblock-chi-homogeneity} 得 \(\chi(x)=\chi_m\) 对所有 \(x\in\mathcal M_m\) 成立。
则成立：
\begin{enumerate}
\normalfont
  \item 最大纤维规模的 \(\ell^\infty\)-型读出：
  \[
  D_m
  =
  \lim_{q\to\infty}\Bigl(2^m\,\tau_{m,1}\!\bigl(\mathrm{IndW}(E_{m,1})^{\,q}\bigr)\Bigr)^{1/(q+1)}.
  \]
  \item \(\chi_m\) 的带符号矩比读出：对任意 \(s,t\in\{0,1\}\) 有
  \[
  \lim_{q\to\infty}
  \frac{\tau_{m,1}\!\bigl(\varepsilon_{\mathrm{scan}}^{\,s}\varepsilon_{\mathrm{rev}}^{\,t}\,\mathrm{IndW}(E_{m,1})^{\,q}\bigr)}
  {\tau_{m,1}\!\bigl(\mathrm{IndW}(E_{m,1})^{\,q}\bigr)}
  =
  \alpha_m^{\,s}\beta_m^{\,t}.
  \]
  \item 扇区内最大纤维规模：对任意 \((\alpha,\beta)\in\{\pm1\}^2\) 定义
  \[
  D_m^{\alpha,\beta}:=\max\{d_m(x):\chi(x)=(\alpha,\beta)\}\qquad(\text{约定空集时为 }0),
  \]
  则
  \[
  D_m^{\alpha,\beta}
  =
  \lim_{q\to\infty}\Bigl(2^m\,\tau_{m,1}\!\bigl(e_{\alpha,\beta}\,\mathrm{IndW}(E_{m,1})^{\,q}\bigr)\Bigr)^{1/(q+1)}.
  \]
\end{enumerate}
\end{corollary}
\begin{proof}
由 \(\widetilde P_x\widetilde P_y=\delta_{xy}\widetilde P_x\) 与 \(\sum_x\widetilde P_x=1\) 得
\[
\mathrm{IndW}(E_{m,1})^{\,q}=\sum_{x\in X_m} d_m(x)^{q}\,\widetilde P_x,
\]
从而
\[
2^m\,\tau_{m,1}\!\bigl(\mathrm{IndW}(E_{m,1})^{\,q}\bigr)
=
\sum_{x\in X_m}d_m(x)^{q+1}.
\]
对有限集合的幂和，\(\lim_{q\to\infty}\bigl(\sum_x d_m(x)^{q+1}\bigr)^{1/(q+1)}=\max_x d_m(x)=D_m\)，得第一条。
\par
第二条由定理 \ref{thm:fold-groupoid-z2x2-joint-spectral-measure} 的带符号矩公式：
\[
2^m\,\tau_{m,1}\!\bigl(\varepsilon_{\mathrm{scan}}^{\,s}\varepsilon_{\mathrm{rev}}^{\,t}\,\mathrm{IndW}(E_{m,1})^{\,q}\bigr)
=
\sum_{x\in X_m}\operatorname{sgn}(c_x)^s\operatorname{sgn}(\iota_x)^t\,d_m(x)^{q+1}.
\]
将分子分母同除以 \(D_m^{q+1}\) 后，所有满足 \(d_m(x)<D_m\) 的项在 \(q\to\infty\) 时消失，仅剩 \(\mathcal M_m\) 的贡献；
而 \(\chi\) 在 \(\mathcal M_m\) 上恒为 \(\chi_m=(\alpha_m,\beta_m)\)，故极限为 \(\alpha_m^s\beta_m^t\)。
\par
第三条由 \(e_{\alpha,\beta}\widetilde P_x=\ind[\chi(x)=(\alpha,\beta)]\,\widetilde P_x\)（定理 \ref{thm:fold-groupoid-z2x2-central-idempotents} 的块投影性质）给出
\[
2^m\,\tau_{m,1}\!\bigl(e_{\alpha,\beta}\,\mathrm{IndW}(E_{m,1})^{\,q}\bigr)
=
\sum_{\chi(x)=(\alpha,\beta)} d_m(x)^{q+1},
\]
再对 \(q\to\infty\) 取 \(\ell^\infty\) 极限即得。证毕。
\end{proof}

\leanverified{paper\_fold\_chi\_layered\_fiber\_spectrum\_certified\_by\_signed\_powersums}
\begin{proposition}[两比特带符号幂和唯一确定 \texorpdfstring{$\chi$}{chi}-分层纤维谱的代数证书]\label{prop:fold-chi-layered-fiber-spectrum-certified-by-signed-powersums}
固定 \(m\)，记 \(N_m:=|X_m|\)。
对 \(s,t\in\{0,1\}\) 与 \(q=0,1,\dots,N_m-1\) 定义两比特带符号幂和
\[
M_{s,t}(q)
:=
2^m\,\tau_{m,1}\!\left(\varepsilon_{\mathrm{scan}}^{\,s}\varepsilon_{\mathrm{rev}}^{\,t}\,\mathrm{IndW}(E_{m,1})^{\,q}\right)
=
\sum_{x\in X_m}\chi_1(x)^s\chi_2(x)^t\,d_m(x)^{q+1},
\]
其中 \(\chi(x)=(\chi_1(x),\chi_2(x))\in\{\pm1\}^2\)。
则对每个 \((\alpha,\beta)\in\{\pm1\}^2\)，多项式
\[
Q_m^{\alpha,\beta}(z):=
\prod_{\substack{x\in X_m\\ \chi(x)=(\alpha,\beta)}}\bigl(1-d_m(x)\,z\bigr)\ \in\ \ZZ[z]
\]
由有限数据族
\(
\{M_{s,t}(q):\ (s,t)\in\{0,1\}^2,\ 0\le q\le N_m-1\}
\)
唯一确定，且该确定性在 Newton 恒等式口径下为可构造的反演。
\end{proposition}
\begin{proof}
对每个 \((\alpha,\beta)\in\{\pm1\}^2\) 与 \(k=1,2,\dots,N_m\)，由 \(\{\pm1\}^2\) 角色正交性（Fourier--Hadamard 反演）得到扇区幂和
\[
p_k^{\alpha,\beta}
:=
\sum_{\chi(x)=(\alpha,\beta)} d_m(x)^k
=
\frac14\sum_{s,t\in\{0,1\}}\alpha^s\beta^t\,M_{s,t}(k-1).
\]
将扇区内的多重集 \(\{d_m(x):\chi(x)=(\alpha,\beta)\}\) 以零根填充到长度 \(N_m\) 后，其幂和仍为 \((p_k^{\alpha,\beta})_{1\le k\le N_m}\)。
Newton 恒等式把有限幂和族 \((p_1,\dots,p_{N_m})\) 与首一多项式
\(
\prod_{j=1}^{N_m}(1-\lambda_j z)
\)
的系数互相唯一确定；
因此 \(Q_m^{\alpha,\beta}(z)\)（零根不贡献因子）由所给数据唯一确定。证毕。
\end{proof}
\leanverified{paper\_fold\_chi\_layered\_fiber\_spectrum\_certified\_by\_signed\_powersums}

\begin{theorem}[window-\texorpdfstring{$6$}{6} 的 \texorpdfstring{$\chi$}{chi}-分层纤维谱可由十二个带符号矩完全重建]\label{thm:fold-window6-chi-layered-spectrum-reconstruction-by-12-moments}
取 \(m=6\)。沿用推论 \ref{cor:fold-bin-m6-fiber-hist} 的记号，则
\[
d_6(x)\in\{2,3,4\},
\qquad
\#\{x\in X_6:\ d_6(x)=2\}=8,\quad
\#\{x\in X_6:\ d_6(x)=3\}=4,\quad
\#\{x\in X_6:\ d_6(x)=4\}=9,
\]
且 \(|X_6|=21\)。再记
\[
\chi(x)=(\chi_1(x),\chi_2(x))\in\{\pm1\}^2,
\]
并对 \(s,t\in\{0,1\}\)、\(q\ge 0\) 定义
\[
M_{s,t}(q)
:=
2^6\,\tau_{6,1}\!\left(
\varepsilon_{\mathrm{scan}}^{\,s}\varepsilon_{\mathrm{rev}}^{\,t}\,
\mathrm{IndW}(E_{6,1})^{\,q}
\right)
=
\sum_{x\in X_6}\chi_1(x)^s\chi_2(x)^t\,d_6(x)^{q+1}.
\]
对每个 \((\alpha,\beta)\in\{\pm1\}^2\) 与 \(j\in\{2,3,4\}\)，记
\[
n_j^{\alpha,\beta}
:=
\#\{x\in X_6:\ \chi(x)=(\alpha,\beta),\ d_6(x)=j\}.
\]
则仅由十二个标量
\[
\{M_{s,t}(q):\ (s,t)\in\{0,1\}^2,\ q=0,1,2\}
\]
便可唯一恢复全部十二个整数 \(n_j^{\alpha,\beta}\)。更精确地，若定义
\[
p_k^{\alpha,\beta}
:=
\sum_{\chi(x)=(\alpha,\beta)} d_6(x)^k
=
\frac14\sum_{s,t\in\{0,1\}}\alpha^s\beta^t\,M_{s,t}(k-1),
\qquad k=1,2,3,
\]
则有显式反演公式
\[
n_2^{\alpha,\beta}
=
3p_1^{\alpha,\beta}-\frac74p_2^{\alpha,\beta}+\frac14p_3^{\alpha,\beta},
\]
\[
n_3^{\alpha,\beta}
=
-\frac83p_1^{\alpha,\beta}+2p_2^{\alpha,\beta}-\frac13p_3^{\alpha,\beta},
\]
\[
n_4^{\alpha,\beta}
=
\frac34p_1^{\alpha,\beta}-\frac58p_2^{\alpha,\beta}+\frac18p_3^{\alpha,\beta}.
\]
从而每个扇区的纤维谱多项式可显式写为
\[
Q_6^{\alpha,\beta}(z)
:=
\prod_{\chi(x)=(\alpha,\beta)}(1-d_6(x)z)
=
(1-2z)^{n_2^{\alpha,\beta}}(1-3z)^{n_3^{\alpha,\beta}}(1-4z)^{n_4^{\alpha,\beta}}.
\]
\end{theorem}
\begin{proof}
由定理 \ref{thm:fold-groupoid-z2x2-joint-spectral-measure} 的带符号矩公式可知，上述 \(M_{s,t}(q)\) 正是 \(\chi\)-分层纤维谱的带符号幂和。再由命题 \ref{prop:fold-chi-layered-fiber-spectrum-certified-by-signed-powersums} 证明中的 Fourier--Hadamard 反演，对每个 \((\alpha,\beta)\in\{\pm1\}^2\) 有
\[
p_k^{\alpha,\beta}
=
\frac14\sum_{s,t\in\{0,1\}}\alpha^s\beta^t\,M_{s,t}(k-1),
\qquad k=1,2,3.
\]

在 \(m=6\) 时，推论 \ref{cor:fold-bin-m6-fiber-hist} 给出全部纤维大小只可能取 \(2,3,4\)。因此对固定扇区 \((\alpha,\beta)\)，必有线性系统
\[
p_1^{\alpha,\beta}=2n_2^{\alpha,\beta}+3n_3^{\alpha,\beta}+4n_4^{\alpha,\beta},
\]
\[
p_2^{\alpha,\beta}=4n_2^{\alpha,\beta}+9n_3^{\alpha,\beta}+16n_4^{\alpha,\beta},
\]
\[
p_3^{\alpha,\beta}=8n_2^{\alpha,\beta}+27n_3^{\alpha,\beta}+64n_4^{\alpha,\beta}.
\]
其系数矩阵
\[
\begin{pmatrix}
2&3&4\\
4&9&16\\
8&27&64
\end{pmatrix}
\]
行列式等于 \(48\neq 0\)，故可逆，直接求逆便得到所示三条反演公式。

最后，由 \(d_6(x)\) 只取 \(2,3,4\) 三个值可知，扇区内多项式 \(Q_6^{\alpha,\beta}(z)\) 的线性因子只能是 \(1-2z\)、\(1-3z\)、\(1-4z\)，其指数分别恰为 \(n_2^{\alpha,\beta}\)、\(n_3^{\alpha,\beta}\)、\(n_4^{\alpha,\beta}\)。故所示因子分解成立。证毕。
\end{proof}
\leanverified{paper\_fold\_window6\_chi\_layered\_spectrum\_reconstruction\_by\_12\_moments}

\begin{proposition}[中心分辨率准则：由 \texorpdfstring{$(\mathrm{IndW},\varepsilon_{\mathrm{scan}},\varepsilon_{\mathrm{rev}})$}{(IndW,eps\_scan,eps\_rev)} 生成的中心子代数的维数与满性判别]\label{prop:fold-center-resolution-criterion-indw-eps}
\leanverified{paper\_fold\_center\_resolution\_criterion\_indw\_eps}
记 \(Z_m:=Z(A_{m,1})\)，并定义中心子代数
\[
Z_m^{(3)}:=\CC\!\left[\mathrm{IndW}(E_{m,1}),\varepsilon_{\mathrm{scan}},\varepsilon_{\mathrm{rev}}\right]\ \subseteq\ Z_m.
\]
则：
\begin{enumerate}
\normalfont
  \item \(Z_m^{(3)}\) 在扇区上分解为
  \[
  Z_m^{(3)}
  =
  \bigoplus_{(\alpha,\beta)\in\{\pm1\}^2}\CC\!\left[\mathrm{IndW}(E_{m,1})\right]\,e_{\alpha,\beta}.
  \]
  \item 令
  \[
  V_{\alpha,\beta}:=\{d_m(x):\ \chi(x)=(\alpha,\beta)\}\subset\NN
  \]
  为扇区内纤维多重度的不同取值集合，则
  \[
  \dim_{\CC}Z_m^{(3)}=\sum_{(\alpha,\beta)\in\{\pm1\}^2}|V_{\alpha,\beta}|.
  \]
  \item 因而 \(Z_m^{(3)}=Z_m\) 当且仅当对每个 \((\alpha,\beta)\) 映射
  \[
  x\longmapsto d_m(x)
  \quad\text{在}\quad
  X_m^{\alpha,\beta}:=\{x\in X_m:\ \chi(x)=(\alpha,\beta)\}
  \quad\text{上为单射。}
  \]
\end{enumerate}
\end{proposition}
\begin{proof}
在 Wedderburn 识别 \(A_{m,1}\cong\bigoplus_{x\in X_m}M_{d_m(x)}(\CC)\) 下，中心 \(Z_m\cong\CC^{X_m}\) 由 \(\{\widetilde P_x\}_{x\in X_m}\) 张成，且
\[
\mathrm{IndW}(E_{m,1})=\sum_x d_m(x)\,\widetilde P_x,\qquad
\varepsilon_{\mathrm{scan}}=\sum_x \chi_1(x)\,\widetilde P_x,\qquad
\varepsilon_{\mathrm{rev}}=\sum_x \chi_2(x)\,\widetilde P_x.
\]
因此任意 \(Z_m^{(3)}\) 元在块 \(x\) 上取值只依赖于 \(\chi(x)\) 与 \(d_m(x)\)，并在同一扇区内为 \(d_m\) 的多项式函数；
而乘以 \(e_{\alpha,\beta}\) 等价于将函数限制到扇区 \(\chi=(\alpha,\beta)\)，得到第一条。
\par
在有限集合上，\(\CC[d]\) 的像空间维数恰为 \(d\) 在该集合上的不同取值数，故第二条成立。
最后由 \(\dim_{\CC}Z_m=|X_m|\) 与第二条立即得到满性判别：当且仅当每个扇区内 \(d_m\) 无碰撞时，\(\sum_{(\alpha,\beta)}|V_{\alpha,\beta}|=|X_m|\)，从而 \(Z_m^{(3)}=Z_m\)。证毕。
\end{proof}

\leanverified{paper\_fold\_center\_resolution\_criterion\_indw\_eps}
\leanverified{paper\_fold\_window6\_center\_three\_observables\_dimension\_defect}
\begin{theorem}[window-\texorpdfstring{$6$}{6} 中三类中心可观测量的维数上限与不可分辨缺陷]\label{thm:fold-window6-center-three-observables-dimension-defect}
取 \(m=6\)，并记
\[
A_{6,1}:=\CC[\mathcal R_6],
\qquad
Z_6:=Z(A_{6,1}),
\]
\[
Z_6^{(3)}
:=
\CC\!\left[
\mathrm{IndW}(E_{6,1}),
\varepsilon_{\mathrm{scan}},
\varepsilon_{\mathrm{rev}}
\right]
\subseteq Z_6.
\]
则必有严格上界
\[
\dim_{\CC}Z_6^{(3)}\le 12.
\]
另一方面，由 \(|X_6|=21\) 得
\[
\dim_{\CC}Z_6=21,
\qquad
\dim_{\CC}(Z_6/Z_6^{(3)})\ge 9.
\]
亦即，仅凭 \(\mathrm{IndW}(E_{6,1})\) 与两比特中心对合 \((\varepsilon_{\mathrm{scan}},\varepsilon_{\mathrm{rev}})\)，不可能完全分辨 window-\(6\) 的全部 \(21\) 个稳定类型；至少有 \(9\) 维中心自由度不会出现在这一三元中心可观测量中。

并且下述条件等价：
\[
\dim_{\CC}Z_6^{(3)}=12
\iff
V_{\alpha,\beta}=\{2,3,4\}
\quad\text{对每个}\quad
(\alpha,\beta)\in\{\pm1\}^2.
\]
\end{theorem}
\begin{proof}
由推论 \ref{cor:fold-bin-m6-fiber-hist}，window-\(6\) 的纤维多重度只取 \(2,3,4\) 三个值，因此对任意 \((\alpha,\beta)\in\{\pm1\}^2\) 都有
\[
V_{\alpha,\beta}\subseteq \{2,3,4\},
\qquad
|V_{\alpha,\beta}|\le 3.
\]
命题 \ref{prop:fold-center-resolution-criterion-indw-eps} 给出维数公式
\[
\dim_{\CC}Z_6^{(3)}
=
\sum_{(\alpha,\beta)\in\{\pm1\}^2}|V_{\alpha,\beta}|,
\]
故立即得到
\[
\dim_{\CC}Z_6^{(3)}\le 4\cdot 3=12.
\]

另一方面，由 Wedderburn 分解
\[
A_{6,1}\cong \bigoplus_{x\in X_6}M_{d_6(x)}(\CC)
\]
可知 \(A_{6,1}\) 恰有 \(|X_6|=21\) 个简单块，从而
\[
\dim_{\CC}Z_6=21.
\]
于是
\[
\dim_{\CC}(Z_6/Z_6^{(3)})
=
\dim_{\CC}Z_6-\dim_{\CC}Z_6^{(3)}
\ge
21-12=9.
\]

最后，\(\dim_{\CC}Z_6^{(3)}=12\) 当且仅当四个扇区都达到上界 \(|V_{\alpha,\beta}|=3\)；而在 \(V_{\alpha,\beta}\subseteq\{2,3,4\}\) 的前提下，这等价于每个扇区都恰好出现全部三种纤维值，即
\[
V_{\alpha,\beta}=\{2,3,4\}
\qquad
(\alpha,\beta\in\{\pm1\}^2).
\]
证毕。
\end{proof}

\begin{conjecture}[扇区内纤维多重度的最终单谱性]\label{conj:fold-sectorwise-eventual-simple-spectrum}
存在 \(m_0\) 使得对所有 \(m\ge m_0\) 及所有 \((\alpha,\beta)\in\{\pm1\}^2\)，映射
\[
d_m:\ X_m^{\alpha,\beta}\longrightarrow \NN
\]
为单射。
若该猜想成立，则由命题 \ref{prop:fold-center-resolution-criterion-indw-eps} 得对所有充分大 \(m\) 有
\[
Z(A_{m,1})=\CC\!\left[\mathrm{IndW}(E_{m,1}),\varepsilon_{\mathrm{scan}},\varepsilon_{\mathrm{rev}}\right],
\]
即中心在制度意义下由“一元谱场 \(\mathrm{IndW}\) 与两比特指纹 \((\varepsilon_{\mathrm{scan}},\varepsilon_{\mathrm{rev}})\)”生成而达到分辨率极限。
\end{conjecture}
